\section{Structure of a simple Ymir program}
\label{sec:basics:program_structure}

Most programs need more than one function. The next listing declares three --
\token{foo}, \token{bar} and \token{main} -- and calls them in a chain. Their
order in the file does not matter to the compiler: a function may call one
declared below it. Custom is nevertheless to put \token{main} last, or in a
module of its own.

\begin{lstlisting}[style=coloredverbatim, caption=First structured source file, label=lst:basics:fst_struct_src]
use std::io;

fn foo() {
  println("Called foo");
  bar(); // calling the bar function
}

fn bar() {
  // because we use std::io, println can be called directly
  println("Called bar");
}

/**
 * The entry point of the program
 */
fn main() {
  println("Program started");
  foo();
}
\end{lstlisting}

The listing also introduces comments. The compiler ignores them entirely; they
exist for whoever reads the code next. A later chapter shows how they can be
turned into documentation automatically. Two syntaxes are available:

\begin{itemize}
\item \token{//} comments out everything to the end of the line.
\item \token{/*} and \token{*/} enclose a comment that may span several lines.
  Starting each of those lines with a star, as on line~14 of
  Listing~\ref{lst:basics:fst_struct_src}, is a convention rather than a rule.
\end{itemize}

\subsection{Source code layout}
\label{sec:basics:source_layout}

The compiler reads a source file as a stream of tokens -- names, parentheses,
commas, brackets and the rest -- and treats any run of whitespace between two of
them as a single separator. Spaces, tabulations and line breaks are
interchangeable, and there may be any number of them. A function can therefore be
written on one line, or spread over a dozen, without changing its meaning.

%% restyle: skip -- spacing here is deliberately arbitrary, that is the point
\begin{lstlisting}[style=coloredverbatim, caption=Arbitrary code layout example]
use std::io;
fn foo(){println("Called foo");bar();}

fn
main
(   )
{
println ("Program started");
foo();
}
\end{lstlisting}

\subsection{Preferred layout}
\label{sec:basics:layout}

That freedom is worth spending on readability rather than on novelty, so the book
follows one layout throughout:

\begin{itemize}
\item one statement per line;
\item two spaces of indentation per level of scope, a scope being what the tokens
  \token{\{\}} delimit;
\item spaces around binary operators;
\item no space before \token{()} or \token{[]}, which attach to the name in front
  of them, as in \token{foo(a, b, c)};
\item no space before the colon that introduces a type, one space after it, as in
  \token{fn foo(a: i32)-> i32}.
\end{itemize}

Every listing in this book is written this way, and Ymir programs are best
written the same way -- uniform code is code a reader can skim.

\vfill
\pagebreak
